\documentclass[11pt]{article}

\usepackage[letterpaper,margin=1in]{geometry}
\usepackage[T1]{fontenc}
\usepackage[utf8]{inputenc}
\usepackage{lmodern}
\usepackage{microtype}
\usepackage{amsmath,amssymb,mathtools,bm}
\usepackage{xcolor}
\usepackage{enumitem}
\usepackage{hyperref}
\usepackage[most]{tcolorbox}
\usepackage{fancyhdr}
\usepackage{parskip}

\definecolor{qftblue}{HTML}{245A8D}
\definecolor{qftgreen}{HTML}{2F6B4F}
\definecolor{qftred}{HTML}{A23B3B}
\definecolor{qftgray}{HTML}{F4F6F8}

\hypersetup{
  colorlinks=true,
  linkcolor=qftblue,
  urlcolor=qftblue,
  pdftitle={PHYS 662 Homework Set 2},
  pdfauthor={Nima Lashkari}
}

\pagestyle{fancy}
\fancyhf{}
\lhead{PHYS 662: Quantum Field Theory I}
\rhead{Homework Set 2}
\cfoot{\thepage}
\setlength{\headheight}{14pt}

\setlist[enumerate]{leftmargin=2.2em,itemsep=0.55em,topsep=0.5em}
\newlist{parts}{enumerate}{2}
\setlist[parts,1]{label=\textbf{(\alph*)},leftmargin=2.3em,itemsep=0.7em}
\setlist[parts,2]{label=\textbf{(\roman*)},leftmargin=2.4em,itemsep=0.45em}

\newcommand{\dd}{\mathrm{d}}
\newcommand{\ii}{\mathrm{i}}
\newcommand{\R}{\mathbb{R}}
\newcommand{\cH}{\mathcal{H}}

\newtcolorbox{conventionbox}[1][]{
  enhanced,
  breakable,
  colback=qftgray,
  colframe=qftblue,
  boxrule=0.8pt,
  arc=1.5mm,
  title={#1},
  fonttitle=\bfseries,
  coltitle=white
}

\newcommand{\problem}[2]{%
  \clearpage
  \section*{\color{qftblue}Problem #1: #2}
  \addcontentsline{toc}{section}{Problem #1: #2}
}


\newcommand{\Lag}{\mathcal L}
\begin{document}

\begin{center}
  {\scshape\Large PHYS 662 --- Quantum Field Theory I\par}
  \vspace{0.25em}
  {\Huge\bfseries\color{qftblue} Homework Set 2\par}
  \vspace{0.25em}
  {\large Lectures 5--6 \quad Fall 2026\par}
\end{center}

\vspace{0.5em}
\noindent
\textbf{Name:} \rule{0.44\textwidth}{0.4pt}
\hfill
\textbf{Due:} \rule{0.24\textwidth}{0.4pt}

\begin{conventionbox}[Instructions and conventions]
Show the reasoning leading to each result. Use \(\hbar=c=k_B=1\) and
\(\eta_{\mu\nu}=\operatorname{diag}(1,-1,-1,-1)\), so
\(\Box=\partial_t^2-\nabla^2\). Work in \(3+1\) dimensions unless a
problem states otherwise. A positive-frequency plane wave is
\(e^{-\ii\omega t+\ii\bm k\cdot\bm x}\).

Keep boundary terms when varying an action, and explain which assumptions
remove them. Field variations vanish at the initial and final times;
spatial boundary conditions are specified below. Treat a complex field and
its conjugate as independent variables when deriving their equations,
then impose complex conjugacy. Derive conservation laws directly when
asked; Noether's first theorem is not required for this set.
\end{conventionbox}

\noindent
The four problems are weighted equally.

\tableofcontents

\problem{1}{The action of a relativistic particle}

Consider a free particle of mass \(m>0\) on a smooth, future-directed
timelike worldline \(x^\mu(\lambda)\) in Minkowski spacetime. Here
\(\lambda\) is an arbitrary increasing parameter, not necessarily time.
Define \(\dot x^\mu=\dd x^\mu/\dd\lambda\) and take
\begin{equation}
 S[x]=-m\int\dd s
 =-m\int_{\lambda_i}^{\lambda_f}\dd\lambda\,
       \sqrt{\eta_{\mu\nu}\dot x^\mu\dot x^\nu},
 \qquad \dot x^2>0.
 \label{eq:particleAction}
\end{equation}
The spacetime endpoints are fixed.

\begin{parts}
 \item Show that the action is invariant under a smooth monotonic
 reparameterization \(\lambda\mapsto\lambda'(\lambda)\) with
 \(\dd\lambda'/\dd\lambda>0\). Explain why reparameterizing the same
 curve does not describe a different physical trajectory.

 \item Vary \(x^\mu(\lambda)\), display the endpoint term, and derive
 the Euler--Lagrange equations before choosing a parameter. Afterward
 set \(\lambda=s\), and show that the trajectory is a straight timelike
 line. Why would imposing \(\dot x^2=1\) in the integrand before taking
 arbitrary variations discard essential information?

 \item Choose coordinate time as the parameter, \(x^0=t\), and write
 the Lagrangian in terms of \(\bm v=\dd\bm x/\dd t\). Calculate
 \begin{equation}
 \bm p=\frac{\partial L}{\partial\bm v},\qquad
 H=\bm p\cdot\bm v-L.
 \end{equation}
 Express \(H\) as a function of \(\bm p\), verify
 \(H^2-|\bm p|^2=m^2\), and derive conservation of \(\bm p\) and
 \(H\) from the equations of motion and the absence of explicit
 position and time dependence.

 \item Expand \(L(\bm v)\) through order \(|\bm v|^4\). Identify the
 rest-energy contribution and the nonrelativistic kinetic term. Explain
 why an additive constant in \(L\) does not change the trajectory but
 changes the reference value of the Hamiltonian.

 \item Explain why simply setting \(m=0\) in
 Eq.~\eqref{eq:particleAction} does not provide a variational principle
 for a massless particle. No alternative massless-particle action is
 required.
\end{parts}

\problem{2}{Maxwell theory in Coulomb gauge}

Consider a prescribed smooth conserved current \(j^\mu=(\rho,\bm j)\),
with \(\partial_t\rho+\nabla\cdot\bm j=0\), and the action
\begin{equation}
 S[A]=\int\dd^4x\left[-\frac14F_{\mu\nu}F^{\mu\nu}
                          -j^\mu A_\mu\right],\qquad
 F_{\mu\nu}=\partial_\mu A_\nu-\partial_\nu A_\mu.
 \label{eq:MaxwellAction}
\end{equation}
Use rationalized units and
\begin{equation}
 A^\mu=(\varphi,\bm A),\qquad
 \bm E=-\nabla\varphi-\partial_t\bm A,\qquad
 \bm B=\nabla\times\bm A.
\end{equation}
The scalar potential \(\varphi(t,\bm x)\) may depend on time. Work on
\(\R^3\) with localized sources and sufficient falloff for the boundary
terms below to vanish; choose potentials vanishing at spatial infinity.

\begin{parts}
 \item Expand the Lagrangian density in terms of \(\varphi\),
 \(\bm A\), \(\rho\), and \(\bm j\). Keep the mixed term involving
 \(\partial_t\bm A\) and \(\nabla\varphi\) at this stage.

 \item Vary \(\varphi\) and all three components of \(\bm A\)
 \emph{before imposing a gauge condition}. Derive their equations and
 identify Gauss's law. Which potential has vanishing canonical momentum,
 and why is its equation a constraint rather than an independent wave
 equation?

 \item Impose Coulomb gauge, \(\nabla\cdot\bm A=0\). Obtain the
 Poisson equation for \(\varphi\) and the sourced equation for
 \(\bm A\). Show that the latter's source is divergence free, using
 current conservation and the Poisson equation. Explain why restricting
 \(\delta\bm A\) to transverse variations from the start would retain
 only the transverse part of its field equation.

 \item Use the spatial Fourier convention
 \begin{equation}
 f(t,\bm x)=\int\frac{\dd^3k}{(2\pi)^3}
     \widetilde f(t,\bm k)e^{\ii\bm k\cdot\bm x}.
 \end{equation}
 For \(k=|\bm k|>0\), define
 \(P_{ij}^{T}(\bm k)=\delta_{ij}-k_i k_j/k^2\) and
 \(\widetilde j_i^{T}=P_{ij}^{T}\widetilde j_j\).
 Find \(\widetilde\varphi\) and solve the forced oscillator equation
 for \(\widetilde{\bm A}\) in terms of \(\widetilde{\bm j}^{T}\) and
 the transverse initial data
 \begin{equation}
 \widetilde{\bm A}(t_0,\bm k)=\bm a(\bm k),\qquad
 \partial_t\widetilde{\bm A}(t_0,\bm k)=\bm b(\bm k),\qquad
 \bm k\cdot\bm a=\bm k\cdot\bm b=0.
 \end{equation}
 Include both the homogeneous solution and the source integral from
 \(t_0\) to \(t\). Spatially constant potentials are excluded by the
 stated falloff conditions.

 \item For a source-free plane wave with \(k>0\), determine the
 dispersion relation and the number of independent polarizations.
 Explain why \(\varphi\) supplies no additional freely propagating
 polarization or independently specifiable initial data.
\end{parts}

\problem{3}{A complex scalar and spontaneous symmetry breaking}

Let \(\Phi\) be a classical complex relativistic scalar with
\begin{equation}
 \Lag=\frac12\partial_\mu\Phi^*\partial^\mu\Phi-V(\Phi),
 \qquad V(\Phi)=\lambda\left(|\Phi|^2-v^2\right)^2,
 \qquad \lambda>0,\quad v>0.
 \label{eq:brokenScalar}
\end{equation}
Our normalization is \(\Phi=\phi_1+\ii\phi_2\), with no factor of
\(1/\sqrt2\); each real component therefore has a canonical kinetic
term \((\partial\phi_a)^2/2\). Use compactly supported variations when
deriving the field equations. Vacuum configurations below refer to
spatially uniform classical fields, with energy compared per unit volume.

\begin{parts}
 \item Derive the equations of motion for \(\Phi\) and \(\Phi^*\).
 Show that the action is invariant under the constant phase rotation
 \(\Phi\mapsto e^{\ii\alpha}\Phi\). Identify the continuous internal
 symmetry group.

 \item Find all constant stationary configurations of the potential,
 classify them, and describe the set of minima. Explain why the vacuum
 minimizes the energy density \(V\), while the Lorentzian action is
 subject to a stationarity condition rather than a general minimization
 principle.

 \item Choose the vacuum \(\Phi_0=v\) and write
 \begin{equation}
 \Phi(t,\bm x)=v+h(t,\bm x)+\ii\chi(t,\bm x),
 \end{equation}
 with real \(h\) and \(\chi\). Expand the Lagrangian through second
 order in these fluctuations. Read off their squared masses from
 canonically normalized quadratic terms, and find both plane-wave
 dispersion relations. Which fluctuation changes the magnitude of the
 order parameter, and which moves tangentially to the vacuum manifold?
 Explain the origin of any massless mode.

 \item Connect this example to Landau's description of phases by
 considering the family
 \begin{equation}
 V_r(\Phi)=\frac r2|\Phi|^2+\lambda|\Phi|^4,
 \qquad \lambda>0.
 \label{eq:LandauFamily}
 \end{equation}
 Find its minima for \(r>0\), \(r=0\), and \(r<0\), and identify
 the value of \(r\) that reproduces Eq.~\eqref{eq:brokenScalar} up to
 an additive constant. Distinguish symmetry of the action from symmetry
 of a chosen vacuum. Explain why internal symmetry restricts the allowed
 potential terms without fixing their coefficients, which can change
 under coarse-graining.

 \item In a regulated finite-volume statistical ensemble with symmetric
 boundary conditions and no external source, explain why
 \(\langle\Phi\rangle=0\) even when the classical potential has a circle
 of minima. Add a weak source \(-\epsilon\operatorname{Re}\Phi\) to the
 energy density. In which order must the limits of infinite spatial
 volume and \(\epsilon\to0^+\) be taken to define a selected ordered
 phase? Assume a regime in which such an ordered phase exists.
\end{parts}

\problem{4}{A nonrelativistic complex field}

In \(d\) spatial dimensions, consider
\begin{equation}
 S[\Psi,\Psi^*]=\int\dd t\,\dd^dx\left[
 \frac{\ii}{2}\left(\Psi^*\partial_t\Psi
                    -(\partial_t\Psi^*)\Psi\right)
 -\frac{1}{2M}\nabla\Psi^*\cdot\nabla\Psi\right],
 \qquad M>0.
 \label{eq:NRAction}
\end{equation}
Here \(\Psi\) is a classical field, with no normalization constraint.
Use periodic spatial boundary conditions in a finite box, or sufficient
falloff on \(\R^d\) so that the integrals and boundary terms are well
behaved.

\begin{parts}
 \item Show that the displayed Lagrangian density is real. Show that
 replacing its time-derivative term by \(\ii\Psi^*\partial_t\Psi\)
 changes the action only by an endpoint term, and state the endpoint
 condition needed for the two forms to give the same bulk equations.
 Why would the candidate kinetic term \(\psi\partial_t\psi\) for a
 single real field fail to produce comparable first-order dynamics?

 \item Vary \(\Psi\) and \(\Psi^*\) independently to obtain both
 equations of motion. For
 \(\Psi=C e^{-\ii\omega t+\ii\bm k\cdot\bm x}\), determine
 \(\omega(\bm k)\). What initial data determine a solution, and how
 do they differ from those of a real Klein--Gordon field?

 \item Compute the canonical momenta
 \(\Pi_{\Psi}=\partial\Lag/\partial(\partial_t\Psi)\) and
 \(\Pi_{\Psi^*}=\partial\Lag/\partial(\partial_t\Psi^*)\).
 Evaluate the energy functional
 \begin{equation}
 H=\int\dd^dx\left[
 \Pi_\Psi\partial_t\Psi+
 \Pi_{\Psi^*}\partial_t\Psi^*-\Lag\right]
 \end{equation}
 and show that it is nonnegative. Explain why the momenta cannot be
 inverted to solve for independent velocities, so that an ordinary
 unconstrained Legendre transform is insufficient. A full constrained
 Hamiltonian analysis is not required.

 \item Starting directly from the equations of motion and their complex
 conjugates, derive a continuity equation for
 \(n=|\Psi|^2\). Find the associated spatial current and state the
 boundary condition under which
 \begin{equation}
 \mathcal N=\int\dd^dx\,|\Psi|^2
 \end{equation}
 is conserved. For a plane wave, compare the ratio of current to density
 with the group velocity \(\nabla_{\bm k}\omega\).

 \item Find the natural mass dimension of \(\Psi\) and check both
 terms in the action dimensionally, using \([M]=1\) and
 \([\partial_t]=[\nabla]=1\). Explain why an equation with the form of
 the one-particle Schrodinger equation does not, by itself, make this
 classical field a normalized one-particle quantum state.
\end{parts}

\end{document}
